\section{Introduction}

Psychiatric diagnoses are clinically indispensable but biologically nonspecific. The
categories of the DSM aggregate patients who differ widely in their biology, course and
treatment response, and they are undercut by pervasive comorbidity, within-category
heterogeneity and diagnostic instability over time~\citep{insel2010rdoc,kotov2017hitop}.
Two influential programmes argue in response that psychopathology is better described by
continuous, \emph{transdiagnostic} dimensions than by categorical boundaries: the U.S.\
National Institute of Mental Health's Research Domain Criteria (RDoC)~\citep{insel2010rdoc}
and the Hierarchical Taxonomy of Psychopathology
(HiTOP)~\citep{kotov2017hitop,ruggero2019hitop}. A recurring motif within this work is a
general factor that spans disorders---the ``\textit{p} factor''---most defensibly read as a
dimension of life impairment rather than a latent liability to specific
symptoms~\citep{caspi2014pfactor,lahey2017general}, though the meaning of such general
factors remains contested~\citep{heinrich2021bifactor,watts2021}.

This dimensional structure has, however, been mapped almost entirely from a single
\emph{kind} of measurement---symptom report~\citep{caspi2014pfactor,kotov2017hitop} or
brain imaging~\citep{marquand2016,wolfers2018,segal2023}. The routinely collected
\emph{biology} that drives the excess physical-health mortality of severe mental
illness---its joint metabolic and inflammatory load---has not been placed \emph{inside} the
transdiagnostic factor space across disorders, even though an independent literature
suggests it forms a partly separable immunometabolic axis that tracks atypical,
energy-related symptoms more than overall
severity~\citep{penninx2024imd,milaneschi2020,lamers2020}, and that metabolic burden
follows medication, adiposity and chronicity more than symptom
load~\citep{vancampfort2015mets,pillinger2017,perry2019}. Whether this load is separable
from a transdiagnostic severity axis, and by how much, has not to our knowledge been
measured directly.

We address this not with a single analysis but with a \textbf{four-stage
measurement-and-validation pipeline}, and we treat the pipeline itself---not any one
finding---as the contribution. Its methodological novelty is not a new estimator but a
\emph{discipline} assembled from three commitments that are rarely held jointly: the map is
estimated from each patient's observed cells only, with \emph{no imputation}, under mixed
likelihoods and a Gaussian-copula continuous block; per-patient \emph{uncertainty is propagated
end-to-end}, so every downstream test (structure, temporal coherence, prognosis) operates on
posterior coordinates rather than point scores; and the map is \emph{fixed
before any validation stage}, with diagnosis held out, so each result is a property of the
measurement rather than an artifact fitted to the outcome. The immunometabolic dissociation that
recurs across all four stages is the existence proof that this discipline reveals structure a
looser measurement would blur. \textbf{(i)~Measurement:} a global, missingness-aware Bayesian
model places routine clinical, cognitive and biological data in one coordinate system with
diagnosis held out, and reports all eight resulting factors uniformly. \textbf{(ii)~Structure:}
on the fixed map, an uncertainty-aware structure test asks whether the patient space is a graded
continuum or a set of discrete kinds, and summarizes it with archetypes rather than imposed
clusters. \textbf{(iii)~Temporal coherence:} the map is scored forward onto follow-up under the
frozen model. \textbf{(iv)~Prognosis:} its incremental value beyond diagnosis and severity is
tested against a hard autoregression bar and benchmarked against the raw indicators.
Two commitments run throughout: we report the eight factors
\emph{uniformly} (loadings, general-factor loadings, factor correlations) and keep
\emph{measurement separate from interpretation}---the clinically salient properties of the map
are read off neutral structural summaries, not built into them. The structural question is itself
contested: normative-modelling studies report individual heterogeneity so extensive that
``the average patient'' is barely informative~\citep{wolfers2018,marquand2016}, whereas the
B-SNIP programme reports replicable, brain-based psychosis ``biotypes'' that cut across DSM
boundaries~\citep{clementz2016biotypes,clementz2020,clementz2024}---a tension whose
resolution may depend on which measurement space is interrogated.

We use latent modelling as a measurement tool, not as a claim that latent factors are more
real than the items that define them. Routine clinical and biological data are
high-dimensional, redundant and pervasively incomplete; a latent model fit to each patient's
observed cells pools many noisy, overlapping indicators into a few axes, places heterogeneous
data types (laboratory values, ordinal scales, counts) on one comparable scale, and gives
every patient a position with quantified uncertainty rather than discarding incomplete
records. Whether that representation \emph{adds} anything over the raw indicators is then an
empirical question, which we test directly: against the full raw indicator set under a matched
classifier, the eight-dimension map is a sufficient representation for predicting deterioration
and near-sufficient for recovery.

Using the harmonized baseline of the three FACE cohorts---bipolar disorder, schizophrenia and
major depression, $N=9{,}013$---we build \emph{FACE-ATLAS}: one global, missingness-aware
Bayesian bifactor / exploratory structural-equation model with mixed likelihoods, fit to each
patient's observed cells only, that places clinical, cognitive, behavioural and laboratory
measurements in one coordinate system while holding diagnosis out as metadata. We then score
the fitted map forward onto two years of follow-up to test whether it persists and whether it
predicts outcome. The pipeline yields an eight-dimension map---a general functional-burden factor
and seven specific axes; a graded continuum rather than discrete biotypes \emph{in this
measurement space}, read through five archetypal extremes; a geometry that persists, with durable
biology and moving symptoms; and a real but modest, \emph{group-level} prognostic signal for
functioning that is co-informative with diagnosis rather than a replacement for it. Running
through all four stages is one recurring empirical property---a single immunometabolic axis that
is distinct from the rest of the clinical picture, the most temporally durable dimension, and the
worst-prognosis pole for functioning---which we present as a property the method \emph{reveals}
rather than the premise it was built on. We report the modest and null results as deliberately as
the positive ones.

Methodologically, FACE-ATLAS proposes a latent-variable dimensional measurement model for a single,
diagnosis-held-out and uncertainty-propagating atlas relying on mixed and Gaussian-copula observation
layers, an observed-cell likelihood for structured missingness, a correlated-\Gfac{} arm that
turns the biology--severity separation into a measured quantity, an archetypal representation
of a continuous patient space, and a downstream errors-in-variables prognosis
analysis \citep{reise2012bifactor,asparouhov2009esem,muthen2012bayesian}. The contribution is the measurement model---its specification and validation
chain---as an \emph{integration discipline} that binds three commitments rarely held jointly
(missingness-native measurement without imputation, uncertainty propagated end-to-end, and a
diagnosis-blind map fixed before validation) rather than as any new generic estimator: routine
biological load is placed inside a transdiagnostic
psychiatric factor space, its separation from functional burden is measured rather than
assumed, and its temporal and prognostic implications are tested without converting a continuum
into biotypes.
